\section{01 Matrix}
\label{sec:graph-01-matrix}

\problemstatement{Given a binary matrix, return a matrix of the same
size where each cell holds the distance to its \textbf{nearest} cell
containing $0$ (distance measured in $4$-directional grid steps).}

\begin{patternbox}
\emph{``Distance to the nearest of several targets''} is
Section~\ref{sec:graph-rotten-oranges}'s multi-source BFS again, applied
to a distance-labeling problem instead of a time-to-completion problem:
seed the queue with every $0$-cell simultaneously (distance $0$), then
BFS outward, where each level reached corresponds to one more unit of
distance.
\end{patternbox}

\subsection*{Algorithm}

\begin{enumerate}
  \item Initialize a $\textit{dist}$ matrix to $0$ wherever the input
        is $0$, and push every such cell into the BFS queue.
  \item For every other cell, initialize $\textit{dist}$ to $-1$
        (unvisited marker) -- this doubles as the \textit{visited}
        check, exactly as a $-1$/sentinel value did in earlier
        chapters.
  \item Run BFS: for each cell popped, examine its $4$ neighbors; any
        neighbor still at $-1$ gets $\textit{dist}[\textit{neighbor}] =
        \textit{dist}[\textit{current}] + 1$ and is pushed.
  \item Return the filled $\textit{dist}$ matrix.
\end{enumerate}

\subsection*{C++ implementation}

\begin{lstlisting}
#include <bits/stdc++.h>
using namespace std;

vector<vector<int>> updateMatrix(vector<vector<int>>& mat) {
    int rows = mat.size(), cols = mat[0].size();
    vector<vector<int>> dist(rows, vector<int>(cols, -1));
    queue<pair<int,int>> q;

    for (int r = 0; r < rows; r++) {
        for (int c = 0; c < cols; c++) {
            if (mat[r][c] == 0) {
                dist[r][c] = 0;
                q.push({r, c});
            }
        }
    }

    int dr[] = {-1, 1, 0, 0};
    int dc[] = {0, 0, -1, 1};

    while (!q.empty()) {
        auto [r, c] = q.front();
        q.pop();

        for (int dir = 0; dir < 4; dir++) {
            int nr = r + dr[dir], nc = c + dc[dir];
            if (nr >= 0 && nr < rows && nc >= 0 && nc < cols &&
                dist[nr][nc] == -1) {
                dist[nr][nc] = dist[r][c] + 1;
                q.push({nr, nc});
            }
        }
    }
    return dist;
}

int main() {
    vector<vector<int>> mat = {
        {0, 0, 0},
        {0, 1, 0},
        {1, 1, 1}
    };
    vector<vector<int>> result = updateMatrix(mat);
    for (auto& row : result) {
        for (int x : row) cout << x << " ";
        cout << endl;
    }
    // 0 0 0
    // 0 1 0
    // 1 2 1
    return 0;
}
\end{lstlisting}

\subsection*{Dry run}

\dryrun{Take the $3\times3$ matrix above.}

Initial queue: every $0$-cell -- $(0,0),(0,1),(0,2),(1,0),(1,2)$, all
at $\textit{dist}=0$. Processing these, their unvisited neighbors get
$\textit{dist}=1$: $(1,1)$ (adjacent to $(0,1)$, $(1,0)$, and $(1,2)$,
whichever processes first), $(2,0)$ (adjacent to $(1,0)$), $(2,2)$
(adjacent to $(1,2)$). Next level: $(2,1)$, adjacent to both
$(1,1)$ (now $\textit{dist}=1$) and $(2,0)$/$(2,2)$ (also
$\textit{dist}=1$), gets $\textit{dist}=2$ (it is reached from a
$\textit{dist}=1$ cell first, since BFS processes level by level).
Final matrix matches the expected output, with the single $\textit{dist}=2$
cell at $(2,1)$ being the only cell two steps from the nearest $0$.

\begin{complexitybox}
\textbf{Time Complexity.} $O(\textit{rows} \times \textit{cols})$ --
each cell is pushed and processed at most once.
\textbf{Space Complexity.} $O(\textit{rows} \times \textit{cols})$ for
the \textit{dist} matrix and the queue.
\end{complexitybox}

\begin{takeawaybox}
Rotten Oranges and 01 Matrix are, structurally, the \emph{same}
algorithm wearing two different costumes: ``minutes until everything is
reached'' and ``distance to the nearest target'' are both just
\textbf{the BFS level number at which each cell is first reached},
read off in two different ways (the level of the very last cell
reached, versus the level of each individual cell respectively).
Recognizing multi-source BFS by its shape -- several simultaneous
starting points, uniform step cost -- matters more than the specific
wording of what the level number represents.
\end{takeawaybox}
